\IfFileExists{stacks-project.cls}{%
\documentclass{stacks-project}
}{%
\documentclass{amsart}
}

% For dealing with references we use the comment environment
\usepackage{verbatim}
\newenvironment{reference}{\comment}{\endcomment}
%\newenvironment{reference}{}{}
\newenvironment{slogan}{\comment}{\endcomment}
\newenvironment{history}{\comment}{\endcomment}

% For commutative diagrams we use Xy-pic
\usepackage[all]{xy}

% We use 2cell for 2-commutative diagrams.
\xyoption{2cell}
\UseAllTwocells

% We use multicol for the list of chapters between chapters
\usepackage{multicol}

% This is generally recommended for better output
\usepackage{lmodern}
\usepackage[T1]{fontenc}

% For cross-file-references
\usepackage{xr-hyper}

% Package for hypertext links:
\usepackage{hyperref}

% For any local file, say "hello.tex" you want to link to please
% use \externaldocument[hello-]{hello}
\externaldocument[introduction-]{introduction}
\externaldocument[conventions-]{conventions}
\externaldocument[sets-]{sets}
\externaldocument[categories-]{categories}
\externaldocument[topology-]{topology}
\externaldocument[sheaves-]{sheaves}
\externaldocument[sites-]{sites}
\externaldocument[stacks-]{stacks}
\externaldocument[fields-]{fields}
\externaldocument[algebra-]{algebra}
\externaldocument[brauer-]{brauer}
\externaldocument[homology-]{homology}
\externaldocument[derived-]{derived}
\externaldocument[simplicial-]{simplicial}
\externaldocument[more-algebra-]{more-algebra}
\externaldocument[smoothing-]{smoothing}
\externaldocument[modules-]{modules}
\externaldocument[sites-modules-]{sites-modules}
\externaldocument[injectives-]{injectives}
\externaldocument[cohomology-]{cohomology}
\externaldocument[sites-cohomology-]{sites-cohomology}
\externaldocument[dga-]{dga}
\externaldocument[dpa-]{dpa}
\externaldocument[sdga-]{sdga}
\externaldocument[hypercovering-]{hypercovering}
\externaldocument[schemes-]{schemes}
\externaldocument[constructions-]{constructions}
\externaldocument[properties-]{properties}
\externaldocument[morphisms-]{morphisms}
\externaldocument[coherent-]{coherent}
\externaldocument[divisors-]{divisors}
\externaldocument[limits-]{limits}
\externaldocument[varieties-]{varieties}
\externaldocument[topologies-]{topologies}
\externaldocument[descent-]{descent}
\externaldocument[perfect-]{perfect}
\externaldocument[more-morphisms-]{more-morphisms}
\externaldocument[flat-]{flat}
\externaldocument[groupoids-]{groupoids}
\externaldocument[more-groupoids-]{more-groupoids}
\externaldocument[etale-]{etale}
\externaldocument[chow-]{chow}
\externaldocument[intersection-]{intersection}
\externaldocument[pic-]{pic}
\externaldocument[weil-]{weil}
\externaldocument[adequate-]{adequate}
\externaldocument[dualizing-]{dualizing}
\externaldocument[duality-]{duality}
\externaldocument[discriminant-]{discriminant}
\externaldocument[derham-]{derham}
\externaldocument[local-cohomology-]{local-cohomology}
\externaldocument[algebraization-]{algebraization}
\externaldocument[curves-]{curves}
\externaldocument[resolve-]{resolve}
\externaldocument[models-]{models}
\externaldocument[functors-]{functors}
\externaldocument[equiv-]{equiv}
\externaldocument[pione-]{pione}
\externaldocument[etale-cohomology-]{etale-cohomology}
\externaldocument[proetale-]{proetale}
\externaldocument[relative-cycles-]{relative-cycles}
\externaldocument[more-etale-]{more-etale}
\externaldocument[trace-]{trace}
\externaldocument[crystalline-]{crystalline}
\externaldocument[spaces-]{spaces}
\externaldocument[spaces-properties-]{spaces-properties}
\externaldocument[spaces-morphisms-]{spaces-morphisms}
\externaldocument[decent-spaces-]{decent-spaces}
\externaldocument[spaces-cohomology-]{spaces-cohomology}
\externaldocument[spaces-limits-]{spaces-limits}
\externaldocument[spaces-divisors-]{spaces-divisors}
\externaldocument[spaces-over-fields-]{spaces-over-fields}
\externaldocument[spaces-topologies-]{spaces-topologies}
\externaldocument[spaces-descent-]{spaces-descent}
\externaldocument[spaces-perfect-]{spaces-perfect}
\externaldocument[spaces-more-morphisms-]{spaces-more-morphisms}
\externaldocument[spaces-flat-]{spaces-flat}
\externaldocument[spaces-groupoids-]{spaces-groupoids}
\externaldocument[spaces-more-groupoids-]{spaces-more-groupoids}
\externaldocument[bootstrap-]{bootstrap}
\externaldocument[spaces-pushouts-]{spaces-pushouts}
\externaldocument[spaces-chow-]{spaces-chow}
\externaldocument[groupoids-quotients-]{groupoids-quotients}
\externaldocument[spaces-more-cohomology-]{spaces-more-cohomology}
\externaldocument[spaces-simplicial-]{spaces-simplicial}
\externaldocument[spaces-duality-]{spaces-duality}
\externaldocument[formal-spaces-]{formal-spaces}
\externaldocument[restricted-]{restricted}
\externaldocument[spaces-resolve-]{spaces-resolve}
\externaldocument[formal-defos-]{formal-defos}
\externaldocument[defos-]{defos}
\externaldocument[cotangent-]{cotangent}
\externaldocument[examples-defos-]{examples-defos}
\externaldocument[algebraic-]{algebraic}
\externaldocument[examples-stacks-]{examples-stacks}
\externaldocument[stacks-sheaves-]{stacks-sheaves}
\externaldocument[criteria-]{criteria}
\externaldocument[artin-]{artin}
\externaldocument[quot-]{quot}
\externaldocument[stacks-properties-]{stacks-properties}
\externaldocument[stacks-morphisms-]{stacks-morphisms}
\externaldocument[stacks-limits-]{stacks-limits}
\externaldocument[stacks-cohomology-]{stacks-cohomology}
\externaldocument[stacks-perfect-]{stacks-perfect}
\externaldocument[stacks-introduction-]{stacks-introduction}
\externaldocument[stacks-more-morphisms-]{stacks-more-morphisms}
\externaldocument[stacks-geometry-]{stacks-geometry}
\externaldocument[moduli-]{moduli}
\externaldocument[moduli-curves-]{moduli-curves}
\externaldocument[examples-]{examples}
\externaldocument[exercises-]{exercises}
\externaldocument[guide-]{guide}
\externaldocument[desirables-]{desirables}
\externaldocument[coding-]{coding}
\externaldocument[obsolete-]{obsolete}
\externaldocument[fdl-]{fdl}
\externaldocument[index-]{index}

% Theorem environments.
%
\theoremstyle{plain}
\newtheorem{theorem}[subsection]{Theorem}
\newtheorem{proposition}[subsection]{Proposition}
\newtheorem{lemma}[subsection]{Lemma}

\theoremstyle{definition}
\newtheorem{definition}[subsection]{Definition}
\newtheorem{example}[subsection]{Example}
\newtheorem{exercise}[subsection]{Exercise}
\newtheorem{situation}[subsection]{Situation}

\theoremstyle{remark}
\newtheorem{remark}[subsection]{Remark}
\newtheorem{remarks}[subsection]{Remarks}

\numberwithin{equation}{subsection}

% Macros
%
\def\lim{\mathop{\mathrm{lim}}\nolimits}
\def\colim{\mathop{\mathrm{colim}}\nolimits}
\def\Spec{\mathop{\mathrm{Spec}}}
\def\Hom{\mathop{\mathrm{Hom}}\nolimits}
\def\Ext{\mathop{\mathrm{Ext}}\nolimits}
\def\SheafHom{\mathop{\mathcal{H}\!\mathit{om}}\nolimits}
\def\SheafExt{\mathop{\mathcal{E}\!\mathit{xt}}\nolimits}
\def\Sch{\mathit{Sch}}
\def\Mor{\mathop{\mathrm{Mor}}\nolimits}
\def\Ob{\mathop{\mathrm{Ob}}\nolimits}
\def\Sh{\mathop{\mathit{Sh}}\nolimits}
\def\NL{\mathop{N\!L}\nolimits}
\def\CH{\mathop{\mathrm{CH}}\nolimits}
\def\proetale{{pro\text{-}\acute{e}tale}}
\def\etale{{\acute{e}tale}}
\def\QCoh{\mathit{QCoh}}
\def\Ker{\mathop{\mathrm{Ker}}}
\def\Im{\mathop{\mathrm{Im}}}
\def\Coker{\mathop{\mathrm{Coker}}}
\def\Coim{\mathop{\mathrm{Coim}}}

% Boxtimes
%
\DeclareMathSymbol{\boxtimes}{\mathbin}{AMSa}{"02}

%
% Macros for moduli stacks/spaces
%
\def\QCohstack{\mathcal{QC}\!\mathit{oh}}
\def\Cohstack{\mathcal{C}\!\mathit{oh}}
\def\Spacesstack{\mathcal{S}\!\mathit{paces}}
\def\Quotfunctor{\mathrm{Quot}}
\def\Hilbfunctor{\mathrm{Hilb}}
\def\Curvesstack{\mathcal{C}\!\mathit{urves}}
\def\Polarizedstack{\mathcal{P}\!\mathit{olarized}}
\def\Complexesstack{\mathcal{C}\!\mathit{omplexes}}
% \Pic is the operator that assigns to X its picard group, usage \Pic(X)
% \Picardstack_{X/B} denotes the Picard stack of X over B
% \Picardfunctor_{X/B} denotes the Picard functor of X over B
\def\Pic{\mathop{\mathrm{Pic}}\nolimits}
\def\Picardstack{\mathcal{P}\!\mathit{ic}}
\def\Picardfunctor{\mathrm{Pic}}
\def\Deformationcategory{\mathcal{D}\!\mathit{ef}}
\usepackage{cleveref}
\crefname{originallemma}{original lemma}{original lemmas}
\Crefname{originallemma}{Original lemma}{Original lemmas}
\makeatletter
\newlabel{lemma-resolution-property-perfect-complex}{{0F8F}{source}{}{sourcepointer-1}{}}
\newlabel{lemma-resolution-property-perfect-complex@cref}{{[originallemma][0][]0F8F}{[1][0][]source}{}{}{}}
\makeatother
\crefname{originallemma}{original lemma}{original lemmas}
\Crefname{originallemma}{Original lemma}{Original lemmas}
\makeatletter
\newlabel{lemma-resolution-property-map-perfect-complex}{{0F8G}{source}{}{sourcepointer-2}{}}
\newlabel{lemma-resolution-property-map-perfect-complex@cref}{{[originallemma][0][]0F8G}{[1][0][]source}{}{}{}}
\makeatother
\crefname{originallemma}{original lemma}{original lemmas}
\Crefname{originallemma}{Original lemma}{Original lemmas}
\makeatletter
\newlabel{lemma-resolution-property-homotopy-map-perfect-complex}{{0F8H}{source}{}{sourcepointer-3}{}}
\newlabel{lemma-resolution-property-homotopy-map-perfect-complex@cref}{{[originallemma][0][]0F8H}{[1][0][]source}{}{}{}}
\makeatother
\crefname{originaldefinition}{original definition}{original definitions}
\Crefname{originaldefinition}{Original definition}{Original definitions}
\makeatletter
\newlabel{derived-definition-cone}{{014E}{source}{}{sourcepointer-4}{}}
\newlabel{derived-definition-cone@cref}{{[originaldefinition][0][]014E}{[1][0][]source}{}{}{}}
\makeatother
\crefname{originalsection}{original section}{original sections}
\Crefname{originalsection}{Original section}{Original sections}
\makeatletter
\newlabel{derived-section-canonical-delta-functor}{{014Z}{source}{}{sourcepointer-5}{}}
\newlabel{derived-section-canonical-delta-functor@cref}{{[originalsection][0][]014Z}{[1][0][]source}{}{}{}}
\makeatother
\crefname{originalsection}{original section}{original sections}
\Crefname{originalsection}{Original section}{Original sections}
\makeatletter
\newlabel{homology-section-truncations}{{0118}{source}{}{sourcepointer-6}{}}
\newlabel{homology-section-truncations@cref}{{[originalsection][0][]0118}{[1][0][]source}{}{}{}}
\makeatother
\title{The K-theory of vector bundles computes the K-theory of perfect complexes on a scheme with the resolution property}
\author{The Stacks Project Authors}
\date{stacks-project:perfect.tex:local-snapshot-20261012; GNU FDL 1.2 or later}
\begin{document}
\maketitle
\setcounter{subsection}{1}
\noindent\textit{Editorial note.} All delivered mathematical text is the author's own, copied from the pinned Stacks Project chapter file perfect.tex; no statement or proof step was rewritten, reordered or removed. The counted span is the author environment and its own proof environment. References inside the span are handed to the assembler as navigation stubs that carry the canonical tag of the cited tag only; whether each of them is load-bearing is recorded on the stub itself.\par\medskip
\noindent\textit{Editorial note.} One sentence inside the counted proof reads, verbatim: 'Details omitted.' (the argument that the constructed map is additive on distinguished triangles, which the author states follows by the same cone argument used a few lines earlier in the proof). The counted span is otherwise the author's complete proof, and the omitted sentence concerns a repetition of an argument that is written out in the same span; per the project rule the phrase is recorded here verbatim.\par\medskip
\noindent\textit{Allowed standard prerequisites (editorial).}\begin{itemize}
\item On a scheme with the resolution property every perfect complex is represented by a bounded complex of finite locally free modules, and the representation is unique up to the usual homotopy arguments. invoked at proof lines 10235-10237 and used throughout. Supplies the representatives on which the alternating sum map is defined and the maps needed for well definedness.
\item A surjection of finite locally free modules is locally split, so its kernel is again finite locally free. used at proof lines 10251-10258. Gives the additivity of the alternating sum on short exact sequences of complexes.
\item An acyclic bounded complex can be filtered by short exact sequences whose terms are the images of the differentials. used at proof lines 10260-10283. Shows that the alternating sum vanishes on acyclic complexes, which is the second ingredient of well definedness.
\item Stupid truncations sit in distinguished triangles whose third term is a single module placed in one degree. invoked at proof lines 10339-10349. Shows by induction that the class of a complex in K-theory is the alternating sum of its terms, giving surjectivity.
\end{itemize}\medskip
\ifdefined\hypertarget\hypertarget{sourcepointer-1}{}\fi
\noindent\textit{Source reference.} Original lemma 0F8F: sources/stacks-perfect/pinned-original-source.tex, source lines 9890--9890. The proof says at lines 10235-10237 that the three resolution property lemmas are used without further mention; they supply the representatives, and every computation with those representatives is inside the counted span.\par\medskip
\ifdefined\hypertarget\hypertarget{sourcepointer-2}{}\fi
\noindent\textit{Source reference.} Original lemma 0F8G: sources/stacks-perfect/pinned-original-source.tex, source lines 9912--9912. One of the three statements used without further mention from proof line 10236; the maps are used inside the counted span for the cone argument.\par\medskip
\ifdefined\hypertarget\hypertarget{sourcepointer-3}{}\fi
\noindent\textit{Source reference.} Original lemma 0F8H: sources/stacks-perfect/pinned-original-source.tex, source lines 9952--9952. The third of the three statements used without further mention from proof line 10236; the compatibility it provides is consumed by the well definedness argument inside the counted span.\par\medskip
\ifdefined\hypertarget\hypertarget{sourcepointer-4}{}\fi
\noindent\textit{Source reference.} Original definition 014E: sources/stacks-chow/tags.txt, source lines 1--1. Cited at proof lines 10303-10305 for the short exact sequence involving the cone; the vanishing used to conclude agreement of the alternating sums is inside the counted span.\par\medskip
\ifdefined\hypertarget\hypertarget{sourcepointer-5}{}\fi
\noindent\textit{Source reference.} Original section 014Z: sources/stacks-chow/tags.txt, source lines 1--1. Cited at proof lines 10307-10309 for the acyclicity of the cone of a quasi-isomorphism; the use made of that acyclicity is inside the counted span.\par\medskip
\ifdefined\hypertarget\hypertarget{sourcepointer-6}{}\fi
\noindent\textit{Source reference.} Original section 0118: sources/stacks-chow/tags.txt, source lines 1--1. Cited at proof lines 10341-10343 for the truncation conventions in the displayed distinguished triangles; the induction that uses them to conclude surjectivity is inside the counted span.\par\medskip
\begin{lemma}
\label{lemma-K-is-old-K}
Let $X$ be a quasi-compact and quasi-separated scheme with the
resolution property. Then the map $K_0(\textit{Vect}(X)) \to K_0(X)$
is an isomorphism.
\end{lemma}

\begin{proof}
This lemma will follow in a straightforward manner from
Lemmas \ref{lemma-resolution-property-perfect-complex},
\ref{lemma-resolution-property-map-perfect-complex}, and
\ref{lemma-resolution-property-homotopy-map-perfect-complex}
whose results we will use without further mention.
Let us construct an inverse map
$$
c : K_0(X) = K_0(D_{perf}(\mathcal{O}_X)) \longrightarrow
K_0(\textit{Vect}(X))
$$
Namely, any object of $D_{perf}(\mathcal{O}_X)$ can be represented
by a bounded complex $\mathcal{E}^\bullet$ of finite locally free
$\mathcal{O}_X$-modules. Then we set
$$
c([\mathcal{E}^\bullet]) = \sum (-1)^i[\mathcal{E}^i]
$$
Of course we have to show that this is well defined. For the moment
we view $c$ as a map defined on bounded complexes of finite locally free
$\mathcal{O}_X$-modules.

\medskip\noindent
Suppose that $\mathcal{E}^\bullet \to \mathcal{F}^\bullet$ is a surjective map
of bounded complexes of finite locally free $\mathcal{O}_X$-modules.
Let $\mathcal{K}^\bullet$ be the kernel. Then we obtain short exact
sequences of $\mathcal{O}_X$-modules
$$
0 \to \mathcal{K}^n \to \mathcal{E}^n \to \mathcal{F}^n \to 0
$$
which are locally split because $\mathcal{F}^n$ is finite locally free.
Hence $\mathcal{K}^\bullet$ is also a bounded complex of finite
locally free $\mathcal{O}_X$-modules and we have
$c(\mathcal{E}^\bullet) = c(\mathcal{K}^\bullet) + c(\mathcal{F}^\bullet)$
in $K_0(\textit{Vect}(X))$.

\medskip\noindent
Suppose given a bounded complex $\mathcal{E}^\bullet$
of finite locally free $\mathcal{O}_X$-modules which is acyclic.
Say $\mathcal{E}^n = 0$ for $n \not \in [a, b]$. Then we
can break $\mathcal{E}^\bullet$ into short exact sequences
$$
\begin{matrix}
0 \to \mathcal{E}^a \to \mathcal{E}^{a + 1} \to \mathcal{F}^{a + 1} \to 0,\\
0 \to \mathcal{F}^{a + 1} \to \mathcal{E}^{a + 2} \to
\mathcal{F}^{a + 3} \to 0, \\
\ldots \\
0 \to \mathcal{F}^{b - 3} \to \mathcal{E}^{b - 2} \to
\mathcal{F}^{b - 2} \to 0, \\
0 \to \mathcal{F}^{b - 2} \to \mathcal{E}^{b - 1} \to \mathcal{E}^b \to 0
\end{matrix}
$$
Arguing by descending induction we see that
$\mathcal{F}^{b - 2}, \ldots, \mathcal{F}^{a + 1}$
are finite locally free $\mathcal{O}_X$-modules, and
$$
c(\mathcal{E}^\bullet) = \sum (-1)[\mathcal{E}^n] =
\sum (-1)^n([\mathcal{F}^{n - 1}] + [\mathcal{F}^n]) = 0
$$
Thus our construction gives zero on acyclic complexes.

\medskip\noindent
It follows from the results of the preceding two paragraphs that $c$
is well defined. Namely, suppose the bounded complexes
$\mathcal{E}^\bullet$ and $\mathcal{F}^\bullet$ of finite locally free
$\mathcal{O}_X$-modules represent the same object of $D(\mathcal{O}_X)$.
Then we can find quasi-isomorphisms
$a : \mathcal{G}^\bullet \to \mathcal{E}^\bullet$ and
$b : \mathcal{G}^\bullet \to \mathcal{F}^\bullet$
with $\mathcal{G}^\bullet$ bounded complex of finite locally free
$\mathcal{O}_X$-modules.
We obtain a short exact sequence of complexes
$$
0 \to \mathcal{E}^\bullet \to C(a)^\bullet \to \mathcal{G}^\bullet[1] \to 0
$$
see Derived Categories, Definition \ref{derived-definition-cone}.
Since $a$ is a quasi-isomorphism, the cone $C(a)^\bullet$ is
acyclic (this follows for example from the discussion in
Derived Categories, Section \ref{derived-section-canonical-delta-functor}).
Hence
$$
0 = c(C(f)^\bullet) = c(\mathcal{E}^\bullet) + c(\mathcal{G}^\bullet[1]) =
c(\mathcal{E}^\bullet) - c(\mathcal{G}^\bullet)
$$
as desired. The same argument using $b$ shows that
$0 = c(\mathcal{F}^\bullet) - c(\mathcal{G}^\bullet)$.
Hence we find that $c(\mathcal{E}^\bullet) = c(\mathcal{F}^\bullet)$
and $c$ is well defined.

\medskip\noindent
A similar argument using the cone on a map
$\mathcal{E}^\bullet \to \mathcal{F}^\bullet$
of bounded complexes of finite locally free $\mathcal{O}_X$-modules
shows that $c(Y) = c(X) + c(Z)$ if $X \to Y \to Z$ is a distinguished triangle
in $D_{perf}(\mathcal{O}_X)$. Details omitted.
Thus we get the desired homomorphism
of abelian groups $c : K_0(X) \to K_0(\textit{Vect}(X))$.

\medskip\noindent
It is clear that the composition
$K_0(\textit{Vect}(X)) \to K_0(X) \to K_0(\textit{Vect}(X))$
is the identity. On the other hand, let $\mathcal{E}^\bullet$
be a bounded complex of finite locally free $\mathcal{O}_X$-modules.
Then the the existence of the distinguished triangles
of ``stupid truncations''
(see Homology, Section \ref{homology-section-truncations})
$$
\sigma_{\geq n}\mathcal{E}^\bullet \to
\sigma_{\geq n - 1}\mathcal{E}^\bullet \to
\mathcal{E}^{n - 1}[-n + 1] \to
(\sigma_{\geq n}\mathcal{E}^\bullet)[1]
$$
and induction show that
$$
[\mathcal{E}^\bullet] = \sum (-1)^i[\mathcal{E}^i[0]]
$$
in $K_0(X) = K_0(D_{perf}(\mathcal{O}_X))$ with apologies for the notation.
Hence the map $K_0(\textit{Vect}(X)) \to K_0(D_{perf}(\mathcal{O}_X)) = K_0(X)$
is surjective which finishes the proof.
\end{proof}

\end{document}
